\documentclass[a4paper,12pt]{article}
\usepackage[margin=2.5cm]{geometry}
\usepackage{hyperref}
\input{macro.tex}

\begin{document}

\begin{center}
{\Large \textbf{Cover Letter}}\\[0.5cm]
{\large Compiler Engineer}\\[0.3cm]
Doyeon Hwang
\end{center}

\vspace{1cm}

% Q1: Motivation
\section{Motivation}

I am applying for a compiler engineer position
based on my experience designing and implementing
a deep learning compiler end-to-end:
from graph analysis and rewrite-pass optimization at the frontend,
to defining a custom MLIR Dialect, building an Emitter,
and writing Linalg lowering patterns at the middle-end.

My core interest lies in transforming deep learning computation graphs
into efficient executable forms
through program analysis and compiler optimization techniques.

Specifically, I built a compiler frontend that parses PyTorch models
via FX graph tracing into a custom IR,
then applies graph-level rewrite passes for optimization.
Beyond the frontend, I constructed an MLIR-based middle-end:
I defined a custom Dialect and its operations using TableGen,
implemented an Emitter that converts frontend IR into MLIR operations,
and wrote OpConversionPattern-based lowering from the custom Dialect to Linalg.
Through this work, I have developed a practical understanding
of the full MLIR pipeline---dialect definition, operation registration,
and multi-stage lowering.

% Q2: What I bring
\section{Relevant Experience and Strengths}

\myparagraph{DL Compiler Frontend Design}
I designed and implemented a compiler frontend
that converts PyTorch models into a custom IR via FX graph tracing.
I defined the Graph, Node, and Value structures from scratch,
clearly separating operation nodes from data flow
and explicitly managing shape, dtype, and producer/consumer relationships
at the IR level.

This design enabled stable application of graph traversal,
cost estimation, and rewrite passes---experience
directly applicable to IR design for analysis and optimization
in compiler infrastructure.

\myparagraph{Graph-level Optimization and Rewrite Passes}
I implemented operator fusion techniques
such as Conv--BatchNorm Folding and Conv--Bias(Add) Folding,
along with Constant Folding and Identity Elimination passes.
Each pass was applied with careful consideration of
preconditions and semantics-preservation constraints,
and I quantitatively verified the effects
by comparing memory access costs before and after optimization.

This process built my ability to reason about
\textit{when} an optimization is safe to apply at the code level---a
skill critical for building reliable model optimization pipelines.

\myparagraph{MLIR-based Middle-end Implementation}
I built a deep learning compiler middle-end using MLIR.
I defined a custom Dialect and operations (Conv, ReLU, Add, etc.)
via TableGen to auto-generate C++ boilerplate,
and implemented an Emitter that converts JSON IR from the frontend
into MLIR operations.
I then wrote OpConversionPattern-based lowering
from the custom Dialect to Linalg,
focusing on correctly expressing each operation's semantics
at the Linalg level---including Conv padding handling
and ReLU-to-linalg::GenericOp conversion.

\myparagraph{Testing, Verification, and IR Design}
During my Ph.D. coursework, I designed and implemented
an automated testing tool for SMT solvers,
detecting 40+ bugs in Z3, CVC5, and bitwuzla.
To unify the I/O formats of heterogeneous solvers,
I designed a custom IR---an experience
that directly informed my IR design decisions
in the deep learning compiler project.

Working through complex system failures,
reducing them to minimal reproducible cases,
and tracing root causes gave me a structured approach
to verifying and debugging system software.

% Q3: Mindset
\section{Problem-solving Approach}

A mentor from industry once told me:

\begin{quote}
\textit{Work on interesting problems, and have ownership.
The most important thing is to find a team that lets you
drive your own problem-solving.}
\end{quote}

This resonated deeply with how I approach engineering.
Rather than treating tasks as requirements to be met,
I tend to dig into \textit{why} a problem exists
and \textit{how far} it can be improved.

Transforming deep learning models into efficient executable forms
is a complex, constraint-heavy problem---but
that complexity is precisely what makes it worth pursuing with ownership.
I want to contribute as an engineer
who takes full responsibility for solving hard problems end-to-end.

% Q4: Research background as strength
\section{Strengths from Research Background}

The core strength of my research background is
the ability to bridge theory and practical system implementation.
I have not only studied compiler and program analysis theory
but also translated those ideas into working code and tools,
always accounting for real-world constraints.

While building the MLIR-based middle-end,
I gained experience rapidly prototyping with AI-assisted tools
and then understanding and refining the core logic myself.
This demonstrates an ability to quickly ramp up on new technology stacks
and apply them to real systems.

Through detecting numerous bugs and tracing their root causes,
I developed a mindset for structurally understanding
and verifying complex systems---an asset
for long-term contributions to model optimization
and compiler verification work.

\vspace{1cm}

\begin{flushright}
Thank you for your consideration.\\[0.5cm]
Doyeon Hwang\\
+82-10-2012-1394\\
fwangdo35@gmail.com
\end{flushright}

\end{document}
